\documentclass[12pt]{beamer}
% \usetheme{umbc4}
\usetheme{moloch}

\usepackage{amsmath,amscd}
\usepackage{fontspec}
\usepackage{unicode-math}
\usepackage[dvipsnames]{xcolor}
\usepackage{nicematrix}


% aet the title bar to black background with white text
\setbeamercolor{frametitle}{bg=black, fg=white}

% ensure other structural elements (bullets, etc.) are visible
\setbeamercolor{structure}{fg=white}

\setmonofont{DejaVu Sans Mono}
\setmainfont{Libertinus Serif}
\setsansfont{Libertinus Sans}
\setmathfont{STIXTwoMath-Regular.otf}
% \setmathfont{STIX Two Math}

% define colors for syntax highlighting
\definecolor{codegreen}{rgb}{0,0.6,0}
\definecolor{codegray}{rgb}{0.5,0.5,0.5}
\definecolor{codecoffee}{rgb}{0.75,0.6,0.3}
\definecolor{question}{rgb}{1.0,0.75,0.1}
\definecolor{lightblue}{rgb}{0.55,0.75,0.9}
\definecolor{lightblue}{rgb}{0.55,0.75,0.9}
\definecolor{yellow}{rgb}{1.0, 1.0, 0.0}
\definecolor{codecomment}{rgb}{0.97, 0.74, 0.4}
\definecolor{codepy}{rgb}{0.9, 0.85, 0.85} %{#F8BC65}
\definecolor{funcColor}{rgb}{0.47, 0.6, 0.8}
\definecolor{midgray}{rgb}{0.4, 0.4, 0.3}
\definecolor{indexwhite}{rgb}{0.9, 0.9, 0.9}

\setbeamercolor{background canvas}{bg=black}
\setbeamercolor{normal text}{fg=white}

\hypersetup{
    colorlinks=true,
    urlcolor=lightblue
}


\newcommand{\FrameTitle}[2]{\frametitle{#1 \hfill \small\textnormal{\textcolor{midgray}{#2}}}}
\newcommand{\ThreePtVspace}{\vspace*{3pt}}
\newcommand{\SixPtVspace}{\vspace*{6pt}}
\newcommand{\nl}{\vspace*{1em}\\}
\newcommand{\vsp}{\vspace*{1em}}

\newcommand{\Matrix}[1]{\textbf{#1}}
\newcommand{\Vector}[1]{\textbf{#1}}
\newcommand{\Scalar}[1]{\textnormal{#1}}
\newcommand{\Quiz}[1]{{\Large\textbf{\color{question}{#1}}}}
\newcommand{\Sketch}[1]{{\Large\textbf{\color{codecoffee}{#1}}}}

\begin{document}
\title{\center{Applied Linear Algebra \\\vsp Linear Functions \\}}

\author{\center{Devi Prasad M \\ Manipal School of Information Sciences \\ Manipal \\}}
\date{\center{August 2026}}
\maketitle

\begin{frame}\FrameTitle{Functions}{Unit 1}
    We write $f: \mathbb{R}^n \rightarrow \mathbb{R}$ to mean that the
    function $f$ maps real $n$-vectors to real numbers.

    \vsp

    Function $f$ is a scalar-valued function of $n$-vectors.

    \vsp
    \textbf{Example}\SixPtVspace
    $$
        f(\Vector{x}:\mathbb{R}^n) = \Vector{1}_n^{\mathrm{T}} \, \Vector{x}
    $$
\end{frame}


\begin{frame}\FrameTitle{Linear Functions}{Unit 1}
    A linear function satisfies the following two properties:\ThreePtVspace
    \begin{itemize}
        \item \emph{Additivity}. $f(\Vector{x} + \Vector{y}) = f(\Vector{x}) + f(\Vector{y})$ for all $n$-vectors $\Vector{x}$ and $\Vector{y}$.
        \vspace*{4pt}
        \item \emph{Homogeneity}. $f(\alpha \Vector{x}) = \alpha\, f(\Vector{x})$ for all $n$-vectors $\Vector{x}$ and all scalars $\alpha$.
    \end{itemize}
\end{frame}

\begin{frame}\FrameTitle{Superposition and Linearity}{Unit 1}
    \vsp
    {\Quiz{Prove the following}}
    \vsp
    \begin{definition}
        \ThreePtVspace
        A function $f$ is \emph{linear} if it satisfies the \emph{superposition principle}:
        \[
            f(\alpha \, \Vector{x} + \beta \, \Vector{y}) = \alpha \, f(\Vector{x}) + \beta \,f(\Vector{y})
        \]
        for all $n$-vectors $\Vector{x}$ and $\Vector{y}$, and all scalars $\alpha$ and $\beta$.
    \end{definition}
\end{frame}

\begin{frame}\FrameTitle{Superposition and Linearity}{Unit 1}
    {\Quiz{Prove the following}}
    \vspace*{4pt}
    \begin{definition}
        \ThreePtVspace
        A function $f$ is \emph{linear} if it satisfies the \emph{superposition principle}:
        \[
            f(\alpha \, \Vector{x} + \beta \, \Vector{y}) = \alpha \, f(\Vector{x}) + \beta \,f(\Vector{y})
        \]
        for all $n$-vectors $\Vector{x}$ and $\Vector{y}$, and all scalars $\alpha$ and $\beta$.
    \end{definition}

    \vsp
    \Sketch{Hint}
    \vspace*{4pt}
    \begin{itemize}
        \item[--] Substitute $\alpha \Vector{x}$ with $\Vector{u}$ and $\beta \Vector{y}$ with $\Vector{v}$.
        \vspace*{4pt}
        \item[--] Apply \emph{addivity} and \emph{homegeneity} to arrive at the result.
    \end{itemize}
\end{frame}

\begin{frame}\FrameTitle{Vector Inner Product Function}{Unit 1}
    Consider the vector inner product function
    \[
        f(\Vector{x}):\mathbb{R} = \Vector{a}^{\textrm{T}} \Vector{x} = a_1x_1 + \cdots + a_nx_n
    \]
    where $\Vector{a}$ is a fixed $n$-vector, and $\Vector{x}$ is a variable $n$-vector.

    \SixPtVspace The function $f$ is called a {\emph{\color{yellow}{linear function}}} of $\Vector{x}$.

    \SixPtVspace The vector $\Vector{a}$ is called the \emph{coefficient vector} of the linear function $f$.

    \SixPtVspace $\Vector{a}$ is also called the \emph{weight vector} of the linear function $f$.
\end{frame}

\begin{frame}\FrameTitle{Vector Inner Product Function}{Unit 1}
    \vsp
    {\Quiz{Prove the following}}\vsp\\

    The vector inner product function
    \[
        f(\Vector{x}):\mathbb{R} = \Vector{a}^{\textrm{T}} \Vector{x} = a_1x_1 + \cdots + a_nx_n, \ \ \  \Vector{a}: \mathbb{R}^n, \ \ \Vector{x}: \mathbb{R}^n
    \]
    is linear.
\end{frame}

\begin{frame}\FrameTitle{Vector Inner Product Function}{Unit 1}
    The vector inner product function
    \[
        f(\Vector{x}):\mathbb{R} = \Vector{a}^{\textrm{T}} \Vector{x} = a_1x_1 + \cdots + a_nx_n, \ \ \  \Vector{a}: \mathbb{R}^n, \ \ \Vector{x}: \mathbb{R}^n
    \]
    is linear.

    \vsp
    \vsp
    {\Quiz{Proof}}\ThreePtVspace\\
    Given the inner product function $f(\Vector{x}) = \Vector{a}^{\textrm{T}} \Vector{x}$, we can verify that it is linear:
    \[
        \begin{array}{rll}
        f(\alpha \, \Vector{x} + \beta \, \Vector{y}) & = & \Vector{a}^{\textrm{T}} (\alpha \, \Vector{x} + \beta \, \Vector{y}) \ThreePtVspace\\
                                      & = & \Vector{a}^{\textrm{T}} (\alpha \, \Vector{x}) + \Vector{a}^{\textrm{T}} (\beta \, \Vector{y}) \ThreePtVspace\\
                                      & = & \alpha (\Vector{a}^{\textrm{T}} \, \Vector{x}) + \beta (\Vector{a}^{\textrm{T}} \, \Vector{y}) \ThreePtVspace\\
                                      & = & \alpha f(\Vector{x}) + \beta f(\Vector{y})
        \end{array}
    \]
\end{frame}

\begin{frame}\FrameTitle{Linear Combination of Vectors}{Unit 1}
    \begin{definition}
        The linear combination of the vectors
        $\Vector{v}_1, \Vector{v}_2, \ldots, \Vector{v}_n$
        with scalars $\Scalar{c}_1, \Scalar{c}_2, \ldots, \Scalar{c}_n$
        is the vector
        \[
            \Scalar{c}_1 \Vector{v}_1 + \Scalar{c}_2 \Vector{v}_2 +
                \cdots + \Scalar{c}_n \Vector{v}_n =
                \sum_{i=1}^{n} \Scalar{c}_{i} \Vector{v}_{i}
        \]
    The scalars are called the \emph{weights} of the linear combination.

    \vsp
    By definition, a linear combination involves only finitely many vectors.
    \end{definition}
\end{frame}

\begin{frame}\FrameTitle{Linear Independece and Span}{Unit 1}
    \begin{definition}[Linear Independence]\ThreePtVspace
         A set of vectors where no single vector can be written
         as a linear combination of the others.
    \end{definition}

    \vsp

    \begin{definition}[Span]\ThreePtVspace
        The complete set of all possible  vectors one can reach by
        changing the scalar values in a linear combination.
    \end{definition}
\end{frame}

\begin{frame}\FrameTitle{Geometric Linearity}{Unit 1}
    \begin{definition}[]
        If a plane $\mathbf{P}$ contains two points
        $A$ and $B$, then it also contains the line $\overrightarrow{AB}$.
    \end{definition}

    \ThreePtVspace
    This expresses the intuitive notion of ``flatness''.

    \ThreePtVspace\ThreePtVspace

    Non-planar surfaces do not satisfy this property.

    \vsp
    \begin{definition}[Geometric Linearity]
        A subset $S$ of $\mathbb{R}^n$ is geometrically linear if $S$
        contains all of the lines through the points of $S$.
    \end{definition}
\end{frame}

\begin{frame}\FrameTitle{Algebraic Linearity}{Unit 1}
    \begin{definition}[]
        A a multivariable polynomial function is called
        algebraically linear if every term has degree 1.
    \end{definition}

    \vsp\vsp
    This is a ``strict'' definition. We may also include
    constants, i.e., terms of degree 0, in which it becomes
    an \emph{affine} function.
\end{frame}

\begin{frame}\FrameTitle{Affine Functions}{Unit 1}
    \begin{definition}[Affine Function]
        An \emph{affine function} is a linear function followed by a
        translation. It can be written as
        \[
            f(\Vector{v}) = \Matrix{L}(\Vector{v}) + \Vector{b}
        \]
        where
        \begin{itemize}
            \item $\Matrix{L}$ is a linear function, and
            \item $\Vector{b}$ is a fixed translation vector.
        \end{itemize}
    \end{definition}
\end{frame}

\begin{frame}\FrameTitle{Affine Function - An Example}{Unit 1}
    Given the general definition of an affine function
    \[
        f(\Vector{v}) = \Matrix{L}(\Vector{v}) + \Vector{b},
    \]
    and the values of the variables,

    \[
        \Matrix{L} =
            \begin{bmatrix}
                2 & 0\\
                0 & 3
            \end{bmatrix},
        \qquad
        \Vector{v} =
            \begin{bmatrix}
            1\\
            2
            \end{bmatrix},
        \qquad
        \Vector{b} =
            \begin{bmatrix}
                4\\
                -1
            \end{bmatrix}.
        \]

        \color{question}{caculate the value of the function $f$.}
\end{frame}


\begin{frame}\FrameTitle{Affine Combination}{Unit 1}
A \emph{linear combination} has the form

\[
    \sum_{i=1}^{k}\Scalar{c}_i\Vector{v}_i,
\]

with no restriction on the coefficients.

\vsp

An \emph{affine combination} has the additional constraint

\[
    \sum_{i=1}^{k}\Scalar{c}_i = 1.
\]

The coefficients may be positive, zero, or negative.
\end{frame}

\begin{frame}\FrameTitle{Convex Combinations}{Unit 1}
In an \emph{affine combination} $\sum_{i=1}^{k}\Scalar{c}_i = 1$,
the coefficients may be positive, zero, or negative.

\vsp\vsp

In addition, if
\[
    \Scalar{c}_i \geq 0 \qquad\text{for all }i,
\]

then the affine combination is called a \emph{convex combination}.
\end{frame}

\begin{frame}\FrameTitle{Regression Model}{Unit 1}
    An $n$-vector \Vector{x} is a \emph{feature vector} of $n$ \emph{features}.\SixPtVspace

    A \emph{regression model} is the affine function of \Vector{x} given by

    \[
        \widehat{\Vector{y}} = \Vector{x}^{\textrm{T}} \beta     + \Vector{v}
    \]

    where $\beta$ is an $n$-vector and \Vector{v} is an $n$-vector.\vsp

    The entries of \Vector{x} are also called the \textbf{regressors} or \textbf{predictors}.\vsp

    \SixPtVspace
        \textbf{regressors} literally mean ``the things that cause the \emph{prediction} or \emph{regression}.''

    \vsp
    $\widehat{\Vector{y}}$ is called the \emph{prediction}.
\end{frame}

\begin{frame}\FrameTitle{Applied Linear Regression}{Unit 1}

Prediction of Intensive Care Unit Length of Stay (LOS)
in the \href{https://www.mdpi.com/2076-3417/13/12/6930}{\emph{MIMIC-IV Dataset}}

\emph{``... four different classification and regression models for predicting LOS in the ICU are
utilized... The first is the common logistic/linear
regression model, which incorporates a linear
combination of all features in modeling...''}

\begin{itemize}
    \item[--] body temperature
    \item[--] chloride
    \item[--] creatinine
    \item[--] heart rate
    \item[--] mean corpuscular volume (MCV)
    \item[--] O2 saturation, and
    \item[--] respiratory rate
\end{itemize}
\end{frame}

\end{document}
